\begin{abstract}
Graph neural networks (GNNs) excel at integrating multimodal patient data for disease prognosis, yet their ``black box'' nature limits clinical adoption. We present a comprehensive explainability framework for the Graph-Integrated Multimodal Attention Network (\giman), addressing interpretability across three Parkinson's disease prediction tasks: diagnostic classification, progression subtype stratification, and prodromal-to-clinical conversion. Our framework combines six complementary methods—attention weight visualization, GNNExplainer subgraph identification, gradient-based feature attribution (IntegratedGradients, GradientSHAP), patient similarity clustering, and counterfactual explanations—applied to a 2,046-patient PPMI cohort. Convergent evidence across methods identifies motor progression rate (UPDRS slope) as the dominant predictor, with 88\% attention coherence validating clinically meaningful patient similarity. Patient clustering achieved 0.45-0.52 silhouette scores, revealing interpretable disease subtypes. Low counterfactual success rates (1/30) confirm prediction robustness. Our framework demonstrates that GNN predictions align with clinical knowledge while providing actionable, patient-level explanations, advancing trustworthy AI in precision medicine.

\noindent\textbf{Keywords:} Graph neural networks, Explainable AI, Parkinson's disease, Precision medicine, Patient similarity networks
\end{abstract}
